% Revision matrix for the proposal defense (task R29).
% 28 Nov 2026: rewritten for the payload router and candidate teacher paths;
% the location column replaces status (an item not yet in the paper says so).
% Section and page numbers are read from build/main.aux, so compile
% main.tex first, then this file:  pdflatex revision-matrix.tex
\documentclass[11pt]{article}
\usepackage[a4paper,margin=2cm,landscape]{geometry}
\usepackage{xr}
\externaldocument{build/main}
\usepackage{longtable,booktabs,array}
\usepackage[T1]{fontenc}
\usepackage{lmodern}
\usepackage{xcolor}
\newcommand{\loc}[1]{\S\ref{#1}, p.~\pageref{#1}}
\newcommand{\notyet}[1]{\textcolor{orange!70!black}{\textit{Not yet in the manuscript:} #1}}
\setlength{\parindent}{0pt}
\begin{document}

{\large\bfseries Revision Matrix}\\[2pt]
\textit{A Risk-Calibrated Adaptive Handover Method for Teacher-Guided Flow-Matching Text-to-Speech Inference}

\medskip
{\small
\setlength{\tabcolsep}{4pt}
\renewcommand{\arraystretch}{1.15}
\begin{longtable}{@{}>{\raggedright\arraybackslash}p{0.6cm} >{\raggedright\arraybackslash}p{2.0cm} >{\raggedright\arraybackslash}p{5.4cm} >{\raggedright\arraybackslash}p{10.2cm} >{\raggedright\arraybackslash}p{5.0cm}@{}}
\toprule
\textbf{No.} & \textbf{Panelist} & \textbf{Comment} & \textbf{How it is addressed} & \textbf{Where in the paper (section, page)} \\
\midrule
\endfirsthead
\toprule
\textbf{No.} & \textbf{Panelist} & \textbf{Comment} & \textbf{How it is addressed} & \textbf{Where in the paper (section, page)} \\
\midrule
\endhead
\midrule\multicolumn{5}{r}{\textit{Continued on next page}}\\
\endfoot
\bottomrule
\endlastfoot
1 & Ms.~Pagal & Are there published studies on the limit or best value? Identify the candidates and features. & No published study gives a best or limiting value; DMOSpeech~2 only suggests ranges. The candidates are three teacher paths of two, four and six larger steps, each cheaper than the released system. The features are the prompt SNR and the invariant residual. & \loc{sec:rrl-handover}; Eq.~\ref{eq:paths}, p.~\pageref{eq:paths}; Table~\ref{tab:cost}, p.~\pageref{tab:cost}; \loc{sec:residual}; \loc{sec:snr} \\
2 & Ms.~Pagal & Is the study looking for the lowest or the highest handover that keeps quality? & For each teacher path, the highest threshold that keeps excess failures within the limit; comparing the calibrated paths then shows the least costly one that stays within it. & \loc{sec:objectives}; Eq.~\ref{eq:objective}, p.~\pageref{eq:objective}; \loc{sec:compared} \\
3 & Ms.~Pagal & State the metrics and the statistical method for comparing the baseline and the method. & Measures: WER, SIM, UTMOS, prosodic variety, failure and excess-failure rates, routing rate, forward passes, saving and runtime. Statistics: paired Wilcoxon tests with Holm correction, exact McNemar tests, equivalence tests, speaker-level bootstrap and exact intervals; the significance level is kept separate from the tolerance $\alpha$. & Table~\ref{tab:measures}, p.~\pageref{tab:measures}; \loc{sec:stats}; \loc{sec:rrl-eval} \\
4 & Ms.~Pagal & What is the basis for counting a candidate as a failure? & One absolute rule for every system: WER above a ceiling, or SIM or UTMOS below minimum levels taken from the reproduced released system on clean development prompts. An excess failure is a failure on an input where the released system does not fail. Failures are counted as a frequency, never discarded. & Definition~\ref{def:failure}, p.~\pageref{def:failure}; Eq.~\ref{eq:excess}, p.~\pageref{eq:excess} \\
5 & Ms.~Pagal & How are the two decision signals compared? & The signals are now the prompt SNR and the invariant residual. The SNR-only router is compared with the full router on every candidate path, next to always and never correcting. The residual is compared with the earlier stabilization signal in the development run. & \loc{sec:compared}; Table~\ref{tab:systems}, p.~\pageref{tab:systems}; Appendix~\ref{app:snr-r}, p.~\pageref{app:snr-r} \\
6 & Dr.~Landicho & Where does the method come from; does the study address an issue or a method? & The method is introduced from the problem: on damaged prompts the teacher's estimate carries the prompt's noise to the student. Each added component is traced to its source, such as histogram equalization from noise-robust speech recognition for the corrected payload. & \loc{sec:background}; Table~\ref{tab:sources}, p.~\pageref{tab:sources} \\
7 & Dr.~Landicho & Justify Learn-Then-Test for TTS; do $\alpha=0.05$ and 90\% confidence come from a study? & LTT needs only pass-or-fail outcomes, independent calibration data and candidates fixed in advance, so it applies to TTS unchanged. The tolerance and confidence are stated as design choices, with a sensitivity check at a larger tolerance. & \loc{sec:rrl-risk}; \loc{sec:calsize} \\
8 & Dr.~Landicho & Lead with the problem; state the gap from the RRL. & The problem statement leads with the fixed handover, its cost and degraded prompts; the gap is stated at the end of the review of literature and of the background. & \loc{sec:problem}; \loc{sec:rrl-gap} \\
9 & Dr.~Landicho & Categorize the related studies (TTS vs.\ speech recognition); distillation vs.\ switching; where the teacher and student are. & The adaptive-computation studies are grouped into five categories and the reviewed studies by domain. The teacher, the student and the handover are placed in the architecture. & Table~\ref{tab:rrl-adaptive}, p.~\pageref{tab:rrl-adaptive}; Table~\ref{tab:synthesis}, p.~\pageref{tab:synthesis}; Fig.~\ref{fig:method}, p.~\pageref{fig:method} \\
10 & Dr.~Landicho & Show that the saving is significant; define the ``cap'' and how cost is measured. & Cost is counted in forward passes, with a guided teacher step counting as two, and runtime on the Tesla T4. The released system's fixed handover (21 passes) replaces the term ``cap''. The worked example gives 10, 13 and 16 passes for the three paths (52\%, 38\% and 24\% fewer). & \loc{sec:baseline}; Table~\ref{tab:cost}, p.~\pageref{tab:cost}; Appendix~\ref{app:calc}, p.~\pageref{app:calc} \\
11 & Dr.~Landicho & Is 0.25 the best threshold? & DMOSpeech~2 uses 0.25 and suggests ranges but does not test whether 0.25 is best; the study keeps the released system as the reference. & \loc{sec:rrl-handover}; \loc{sec:scope} \\
12 & Dr.~Landicho & Calibration is part of the method, not a problem; all problems are anchored in the handover. & One general problem, anchored in the cost of the fixed handover; calibration is part of Objective~2, which develops and calibrates the router. & \loc{sec:problem}; \loc{sec:objectives}; \loc{sec:obj2}; \loc{sec:obj3} \\
13 & Dr.~Landicho & Consider a pre-experiment to validate the problem. & The preliminary experiments show that a fixed early handover fails on noisy prompts and that adapting the handover time could not be certified; they motivate the design. & Appendix~\ref{app:prelim}, p.~\pageref{app:prelim}; \loc{sec:architecture} \\
14 & Dr.~Landicho & Do not set the baseline failure to zero; measure DMOSpeech~2's quality by reproducing it. & The released system is judged by the same failure rule as every other system, and its scores are reproduced and compared with the published values before use. & Definition~\ref{def:failure}, p.~\pageref{def:failure}; \loc{sec:obj1} \\
15 & Dr.~Landicho & Keep the lightweight goal: low computation with good quality. & Goal stated as a fixed, lower cost per utterance with excess failures within a stated limit; experiments on the Tesla T4; on-device use outside the scope. & \loc{sec:background}; \loc{sec:scope} \\
16 & Dr.~Landicho & Establish a threshold future researchers can use. & The calibrated threshold of each teacher path is reported with the conditions under which it holds and a procedure for recalibrating it. & \loc{sec:verify}; \loc{sec:significance} \\
17 & Dr.~Landicho & Establish the RRL; define the metrics and how the comparison is made. & The metrics, their limits and the paired comparison are stated in the review of literature and in the evaluation plan. & \loc{sec:rrl-eval}; \loc{sec:compared}; \loc{sec:stats} \\
18 & Dr.~Landicho & Establish novelty against flow-matching work in image generation. & Image and video methods adapt how much computation is spent or when the handover happens; none adapts what is handed over, attaches a statistical guarantee to it, or addresses speech synthesis. The sources of the added components are listed separately. & \loc{sec:background}; \loc{sec:rrl-adaptive}; Table~\ref{tab:sources}, p.~\pageref{tab:sources} \\
19 & Dr.~Landicho & Use moderate noise; test on an unseen dataset. & Three moderate conditions at 25~dB (white, pink and high-pass noise) are among the eight prompt conditions; LibriSpeech-PC test-clean is used only for testing. & \loc{sec:data}; Table~\ref{tab:conditions}, p.~\pageref{tab:conditions} \\
20 & Dr.~Landicho & The experiment does not seem complex enough; look into other problems as well.$^{\ast}$ & The evaluation covers eight prompt conditions, three calibrated teacher paths, four ablations, a hindsight bound, a pre-specified test of prosodic variety with a listening check, a single and a repeated check of the guarantee, and robustness to the prompt condition, calibration on clean prompts, smaller calibration groups, a larger tolerance and a second dataset. & \loc{sec:compared}; \loc{sec:verify}; \loc{sec:robust}; \loc{sec:stats} \\
21 & Ms.~Eder & Create a comparison matrix against the baseline; this is where the contribution is shown. & Every compared system will be reported on every measure, on the same inputs, overall, per prompt condition and on the second dataset. & Table~\ref{tab:systems}, p.~\pageref{tab:systems}; Table~\ref{tab:measures}, p.~\pageref{tab:measures} \\
22 & Ms.~Eder & Setting up the threshold is the main contribution; back it with literature. & Calibrated thresholds of this kind are cited from CALM, thought calibration, FlowGuard, adaptive ASR correction and selective generative enhancement; the method calibrates a threshold for each teacher path with a statistical guarantee. & \loc{sec:background}; \loc{sec:rrl-adaptive}; \loc{sec:rrl-risk}; \loc{sec:calibration} \\
\end{longtable}
}
{\footnotesize $^{\ast}$From the transcript (Recording~26, 02:24); not in the minutes.\par}
\end{document}
